\chapter{System Architecture}

\section{Module Organization}

The system is organized into distinct modules with clear responsibilities and minimal coupling.

\subsection{Core Modules}

\subsubsection{src/core/config.py}

\textbf{Purpose:} Centralized configuration management

\textbf{Key Components:}
\begin{itemize}
    \item Project root path determination
    \item Data directory paths
    \item Default parameters
\end{itemize}

\begin{lstlisting}[style=pythonstyle, caption=Configuration Module Structure]
from pathlib import Path

# Determine project root
PROJECT_ROOT = Path(__file__).parent.parent.parent

# Data paths
DATA_RAW = PROJECT_ROOT / "data" / "raw"
DATA_OUTPUTS = PROJECT_ROOT / "outputs"

# Default parameters
DEFAULT_ID_ROW = 7
DEFAULT_CODE_COL = 0
DEFAULT_LABEL_COL = 2
\end{lstlisting}

\subsubsection{src/core/normalization.py}

\textbf{Purpose:} String normalization for consistent matching

\textbf{Key Functions:}
\begin{itemize}
    \item \texttt{\_norm\_label(s: str) -> str}: Normalize column labels
    \item \texttt{\_norm\_emp\_id(s: str) -> str}: Normalize employee IDs
\end{itemize}

\textbf{Normalization Rules:}
\begin{enumerate}
    \item Convert to lowercase
    \item Remove accents and diacritics
    \item Strip whitespace
    \item Remove special characters
    \item Replace hyphens/underscores with spaces
\end{enumerate}

\begin{lstlisting}[style=pythonstyle, caption=Normalization Function]
import unicodedata
import re

def _norm_label(s: str) -> str:
    """Normalize label for consistent matching."""
    if not s:
        return ""
    # Convert to string and lowercase
    s = str(s).lower().strip()
    # Remove accents
    s = unicodedata.normalize('NFKD', s)
    s = s.encode('ascii', 'ignore').decode('ascii')
    # Remove special chars
    s = re.sub(r'[^a-z0-9\s]', '', s)
    # Collapse whitespace
    s = re.sub(r'\s+', ' ', s).strip()
    return s
\end{lstlisting}

\textbf{Why Normalization Matters:}

Excel files from different sources may have inconsistent labeling:
\begin{itemize}
    \item "Salaire Brut" vs "salaire\_brut" vs "SALAIRE BRUT"
    \item "Cotisation Salariale" vs "cot. salariale"
    \item Employee IDs: "EMP-001" vs "EMP001" vs "emp 001"
\end{itemize}

Normalization ensures these variants match correctly.

\subsection{Data Loading Modules}

\subsubsection{src/core/multi\_sheet\_loader.py}

\textbf{Purpose:} Load synthetic payroll data from Excel multi-sheet structure

\textbf{Key Function:}
\begin{lstlisting}[style=pythonstyle]
def load_all_months(
    workbook_path: str,
    output_csv: str,
    month_regex: str = r"^\d{4}-\d{2}$"
) -> pd.DataFrame:
    """
    Load all month sheets from workbook.
    
    Args:
        workbook_path: Path to Excel workbook
        output_csv: Path to save consolidated CSV
        month_regex: Regex pattern for month sheet names
        
    Returns:
        DataFrame with all employee-month records
    """
\end{lstlisting}

\textbf{Processing Steps:}
\begin{enumerate}
    \item Open workbook with xlwings
    \item Identify month sheets by regex pattern
    \item For each month sheet:
    \begin{itemize}
        \item Read header row
        \item Read employee data
        \item Normalize column names
        \item Add month column
    \end{itemize}
    \item Concatenate all months
    \item Standardize column order
    \item Save to CSV
\end{enumerate}

\subsubsection{tests/mapped\_data\_loader.py}

\textbf{Purpose:} Load real payroll data after mapping

\textbf{Key Differences from multi\_sheet\_loader:}
\begin{itemize}
    \item Handles MM format month names (converts to YYYY-MM)
    \item Inserts role column (for compatibility)
    \item Reorders columns to match synthetic data
    \item More defensive error handling
\end{itemize}

\textbf{Month Format Conversion:}
\begin{lstlisting}[style=pythonstyle]
# Detect format
if df["month"].str.match(r"^\d{2}$").all():
    # Convert MM to YYYY-MM with base year
    month_strings = df["month"].astype(str).str.zfill(2)
    df["month"] = "2025-" + month_strings
    print("Converted MM format to YYYY-MM (base year 2025)")
\end{lstlisting}

\subsection{Excel Integration Modules}

\subsubsection{src/excel\_tools/mapping\_utils.py}

\textbf{Purpose:} Shared utilities for mapping operations

\textbf{Key Functions:}

\begin{table}[H]
\centering
\begin{tabularx}{\textwidth}{lX}
\toprule
\textbf{Function} & \textbf{Purpose} \\
\midrule
\texttt{ensure\_custom\_sheet} & Create or retrieve CustomMap sheet \\
\texttt{serialize\_keys\_ops} & Convert formula list to storage format \\
\texttt{get\_source\_rows\_display} & Extract displayable source rows \\
\texttt{save\_mapping\_row} & Save/update mapping in CustomMap \\
\texttt{get\_month\_sheets} & Identify month sheets by regex \\
\bottomrule
\end{tabularx}
\caption{Mapping Utility Functions}
\end{table}

\textbf{Storage Format:}

Mappings are stored in CustomMap sheet as:
\begin{lstlisting}
Destination_Label | Mapping_Type | KeysOps
salaire_brut      | formula      | +::SALAIRE_BASE::label;;+::PRIME::label
\end{lstlisting}

This serialized format allows:
\begin{itemize}
    \item Persistence across sessions
    \item Easy parsing and reconstruction
    \item Support for complex multi-term formulas
\end{itemize}

\subsubsection{src/excel\_tools/import\_dialog.py}

\textbf{Purpose:} Main entry point for Excel import workflow

\textbf{Responsibilities:}
\begin{enumerate}
    \item File selection dialog
    \item Launch mapping GUI
    \item Apply mappings automatically
    \item Generate CSV
    \item Handle errors gracefully
\end{enumerate}

\textbf{Workflow:}

\begin{figure}[H]
\centering
\begin{tikzpicture}[
    node distance=1.2cm,
    startstop/.style={rectangle, rounded corners, draw, fill=red!20, text width=3cm, align=center, minimum height=1cm},
    process/.style={rectangle, draw, fill=blue!20, text width=3.5cm, align=center, minimum height=1cm},
    decision/.style={diamond, draw, fill=green!20, text width=2.5cm, align=center, aspect=2},
    arrow/.style={->, >=stealth, thick}
]

\node[startstop] (start) {Start};
\node[process, below of=start] (load) {Load data/raw/*.xlsx files};
\node[process, below of=load] (display) {Display file list};
\node[decision, below of=display] (select) {File selected?};
\node[process, below of=select, yshift=-0.5cm] (check) {Check data\_sorted.xlsx};
\node[decision, below of=check] (exists) {Exists?};
\node[process, left of=exists, xshift=-2cm] (create) {Create new workbook};
\node[process, right of=exists, xshift=2cm] (open) {Open existing};
\node[process, below of=exists, yshift=-1cm] (gui) {Launch mapping GUI};
\node[process, below of=gui] (apply) {Apply mappings};
\node[process, below of=apply] (save) {Save to data\_sorted.xlsx};
\node[process, below of=save] (csv) {Generate CSV};
\node[startstop, below of=csv] (end) {End};

\draw[arrow] (start) -- (load);
\draw[arrow] (load) -- (display);
\draw[arrow] (display) -- (select);
\draw[arrow] (select) -- node[right] {Yes} (check);
\draw[arrow] (select) -- node[above] {No} +(-3,0) |- (display);
\draw[arrow] (check) -- (exists);
\draw[arrow] (exists) -- node[above] {No} (create);
\draw[arrow] (exists) -- node[above] {Yes} (open);
\draw[arrow] (create) |- (gui);
\draw[arrow] (open) |- (gui);
\draw[arrow] (gui) -- (apply);
\draw[arrow] (apply) -- (save);
\draw[arrow] (save) -- (csv);
\draw[arrow] (csv) -- (end);

\end{tikzpicture}
\caption{Import Dialog Workflow}
\end{figure}

\subsection{Mapping Engine}

\subsubsection{tests/test\_monthly\_mapping\_engine.py}

\textbf{Purpose:} Core mapping calculation engine

\textbf{Key Components:}

\begin{algorithm}[H]
\caption{Mapping Application Algorithm}
\begin{algorithmic}[1]
\State \textbf{Input:} Source sheet, Destination sheet, Mapping rules
\State \textbf{Output:} Populated destination sheet

\For{each month sheet}
    \State Detect employee blocks on source
    \State Build row lookup (code/label $\rightarrow$ row number)
    \For{each employee}
        \State Copy employee ID to destination
        \For{each mapping rule}
            \State Parse formula (operators and keys)
            \State Initialize result = 0
            \For{each term in formula}
                \State Find source row by key
                \State Extract value from employee's column
                \State Apply operator (+, -, *, /)
                \State Update result
            \EndFor
            \State Write result to destination cell
        \EndFor
    \EndFor
\EndFor
\end{algorithmic}
\end{algorithm}

\textbf{Formula Parsing:}

Example: \texttt{salaire\_brut = +SALAIRE\_BASE +PRIME -AVANCE}

\begin{enumerate}
    \item Split by operators: \texttt{["", "SALAIRE\_BASE", "PRIME", "AVANCE"]}
    \item Extract operators: \texttt{["+", "+", "-"]}
    \item For each term:
    \begin{itemize}
        \item Lookup row: SALAIRE\_BASE $\rightarrow$ row 5
        \item Get employee column: Employee 1 $\rightarrow$ col C
        \item Read value: sheet.range("C5").value $\rightarrow$ 5000
        \item Apply operator: result + 5000
    \end{itemize}
\end{enumerate}

\subsection{Analytics Modules}

\subsubsection{src/features/analytics/dashboard\_gui.py}

\textbf{Purpose:} Interactive visualization dashboard

\textbf{Architecture:}

\begin{lstlisting}[style=pythonstyle, caption=Dashboard Class Structure]
class PayrollDashboard:
    def __init__(self):
        self.root = tk.Tk()
        self.df = None  # Payroll data
        self.setup_ui()
        self.load_data()
    
    def setup_ui(self):
        """Create GUI layout with buttons."""
        # Left panel: plot buttons
        # Right panel: matplotlib canvas
    
    def load_data(self):
        """Load CSV and validate."""
    
    def plot_total_cost_trend(self):
        """Plot 1: Total cost over time."""
    
    # ... 9 more plot methods
    
    def run(self):
        """Start GUI main loop."""
        self.root.mainloop()
\end{lstlisting}

\textbf{Plot Generation Pattern:}

Each plot method follows this pattern:
\begin{enumerate}
    \item Validate data loaded
    \item Clear previous plot
    \item Process data (group, filter, aggregate)
    \item Create matplotlib figure
    \item Customize appearance (colors, labels, grid)
    \item Embed in tkinter canvas
\end{enumerate}

\section{Data Flow Architecture}

\subsection{Real Data Pipeline}

\begin{figure}[H]
\centering
\begin{tikzpicture}[
    node distance=1cm,
    box/.style={rectangle, draw, fill=blue!15, text width=3.5cm, align=center, minimum height=0.8cm, font=\small},
    arrow/.style={->, >=stealth, thick}
]

\node[box] (excel) {Excel File \\ (essai donnees paye.xlsx)};
\node[box, below=of excel] (import) {import\_dialog.py \\ File Selection};
\node[box, below=of import] (gui) {Mapping GUI \\ Formula Builder};
\node[box, below=of gui] (custom) {CustomMap Sheet \\ Persistent Storage};
\node[box, below=of custom] (engine) {Mapping Engine \\ Calculation};
\node[box, below=of engine] (sorted) {data\_sorted.xlsx \\ Mapped Data};
\node[box, below=of sorted] (loader) {mapped\_data\_loader.py \\ CSV Conversion};
\node[box, below=of loader] (csv) {payroll\_mapped\_long.csv};

\draw[arrow] (excel) -- (import);
\draw[arrow] (import) -- (gui);
\draw[arrow] (gui) -- (custom);
\draw[arrow] (custom) -- (engine);
\draw[arrow] (engine) -- (sorted);
\draw[arrow] (sorted) -- (loader);
\draw[arrow] (loader) -- (csv);

\end{tikzpicture}
\caption{Real Data Processing Pipeline}
\end{figure}

\subsection{Synthetic Data Pipeline}

\begin{figure}[H]
\centering
\begin{tikzpicture}[
    node distance=1cm,
    box/.style={rectangle, draw, fill=green!15, text width=3.5cm, align=center, minimum height=0.8cm, font=\small},
    arrow/.style={->, >=stealth, thick}
]

\node[box] (gen) {generate\_synthetic\_payroll.py \\ Data Generation};
\node[box, below=of gen] (excel) {Synthetic Workbook \\ Multi-sheet Structure};
\node[box, below=of excel] (loader) {multi\_sheet\_loader.py \\ Direct CSV Conversion};
\node[box, below=of loader] (csv) {payroll\_long.csv};

\draw[arrow] (gen) -- (excel);
\draw[arrow] (excel) -- (loader);
\draw[arrow] (loader) -- (csv);

\end{tikzpicture}
\caption{Synthetic Data Processing Pipeline}
\end{figure}

\subsection{Unified Analytics}

Both pipelines produce identical CSV formats, enabling:

\begin{lstlisting}[style=pythonstyle]
# Load either dataset
df_real = pd.read_csv("outputs/payroll_mapped_long.csv")
df_synthetic = pd.read_csv("outputs/payroll_long.csv")

# Same columns, same format
assert df_real.columns.equals(df_synthetic.columns)

# Same analysis code works on both
def analyze(df):
    monthly = df.groupby('month')['total_cost'].sum()
    return monthly.mean(), monthly.std()

mean_real, std_real = analyze(df_real)
mean_synth, std_synth = analyze(df_synthetic)
\end{lstlisting}

This interchangeability allows:
\begin{itemize}
    \item Validate methodology on real data
    \item Run actual analysis on synthetic data
    \item Compare results between datasets
    \item Switch datasets without code changes
\end{itemize}
